\cotefile{71}
\scanpage{43}\mspage{XV}
Choisissons un voisinage $V$ de $\Gamma_{x'y'}$ tel que : $\mathrm{I})\ V\cap\Gamma_{ss'}=\unsure{\phi|\varnothing}$ ;
% Paired vertical strokes at the two edges of the upper sheet.
\reversemarginpar\margin{$\vert$}\normalmarginpar\margin{$\vert$}%
$\overline{\underline{\mathrm{II}}})\ V\cap\mathrm{int}\,\Lambda^{2}$, $V\cap\mathrm{ext}\,\Lambda^{2}$ sont des ouverts connexes. (Il \unsure{\.{y}|\'{y}|y} en a toujours car si $A$ et $B$ sont deux parties compactes de $\mathbb{R}^{2}$ il \unsure{exìste|existe} un couple de voisinages $A\subset V_A$, $B\subset V_B$ tel que $V_A\cap V_B=\phi$, le lemme I \deleted{\unsure{fai|fait}} montre qu'on peut trouver $V\inserted{\subset V_A}$ satisfaisant en plus $\overline{\underline{\mathrm{II}}})$ si $A=\Gamma_{\deleted{\unsure{s}}x'\deleted{\unsure{s}}y'}$ et $B=\Gamma_{ss'}$)

Comme $\Lambda_{xx'}\cap V\inserted{-\{x'\}}\ne\phi$ \deleted{\illegible{1}} (resp. $\Lambda_{y'y}\cap V-\unsure{\{v'\}|\{y'\}}\ne\phi$) on a
\[
\begin{aligned}
&\Lambda_{xx'}\cap V-\underbrace{\Lambda_{xx'}\cap\Lambda^{2}}_{x'}\subset V\cap\mathrm{int}\,\Lambda^{2}\ \cup\ V\cap\mathrm{ext}\,\Lambda^{2}\quad(\text{resp.}\\
&\Lambda_{y'y}\cap V-\underbrace{\Lambda_{y'y}\cap\Lambda^{2}}_{y'}\subset V\cap\mathrm{int}\,\Lambda^{2}\ \cup\ V\cap\mathrm{ext}\,\Lambda^{2})\quad\text{mais alors}\\
&\Lambda_{xx'}\cap V-\{x'\}\subset V\cap\mathrm{ext}\,\Lambda^{2}\quad(\text{resp. }\Lambda_{y'y}\cap V-\{y'\}\subset V\cap\mathrm{ext}\,\Lambda^{2})\\
&\text{car}\quad \Lambda_{xx'}-\{x'\}\subset\mathrm{ext}\,\Lambda^{2}\quad(\text{resp. }\Lambda_{y'y}-\{y'\}\subset\mathrm{ext}\,\Lambda^{2}).
\end{aligned}
\]
Soit donc $x'',\deleted{y''\in}\ \in(\Lambda_{xx'}-\{x'\})\cap V$ (resp. $y''\in(\Lambda_{y'y}\unsure{\leftarrow|-}\{y'\})\cap V$).
Il est clair qu'on peut rélier $x''$ à $y''$ par un arc de polygone $\inserted{\Lambda_{x''y''}}\subset V$ avec $\Lambda_{x''y''}\cap\Gamma^{\unsure{\prime|1}}_{ss'}=\phi$ (cela résulte \unsure{dee|deu|du} fait que $V\cap\Gamma'_{ss'}=\phi$) alors l'arc de polygone $\inserted{\overline{\Lambda}_{xy}}$ obtenu
\reversemarginpar\margin{$\vert$}\normalmarginpar\margin{$\vert$}%
en récollant $\Lambda_{xx''}$ avec $\Lambda_{x''y''}$, $\Lambda_{x''y''}$ avec $\Lambda_{y''y}$, selon $\{x''\}$ et $\{y''\}$ respectivement, satisfait $\overline{\Lambda}_{x\unsure{\dot y|y}}\subset U=\unsure{\mathrm{Int}|\mathrm{int}}\,\Lambda^{1}-\unsure{R|\mathbb{R}}_{\Lambda^{2}}\cap\mathrm{int}\,\Lambda^{1}$.

Prouvons que $V$ est connexe. Soient $x,y\in U=\mathrm{ext}\,\Lambda^{1}\cup\mathrm{int}\,\Lambda^{2}\cup{]s,s'[}$

Si $x,y\in\mathrm{ext}\,\Lambda^{1}$ on rélie $x$ à $y$ par un arc de polygone $\Lambda_{xy}\subset\mathrm{ext}\,\Lambda^{1}$.

Si $x,y\in\mathrm{int}\,\Lambda^{2}\cup{]s,s'[}$ on peut aussi rélier $x$ à $y$ par un arc de polygone $\Lambda_{xy}\subset\mathrm{ext}\,\Lambda^{1}$.

Donc il reste le cas ou $x\unsure{\in|\notin}\mathrm{ext}\,\Lambda^{1}$ et $y\in\mathrm{int}\,\Lambda^{2}\cup{]s,s'[}$. On choisit $z\in{]s,s'[}$ on rélie $x$ a $z$ par $\Lambda_{xz}$ tel que $\Lambda_{xz}-\{z\}\subset\mathrm{ext}\,\Lambda^{1}$ et $y$ à $z$ par $\Lambda_{zy}$ tel que $\Lambda_{zy}\deleted{\unsure{-\{z\}}}\subset\mathrm{int}\,\Lambda^{2}\cup{]s,s'[}$. Il est \deleted{\unsure{q}}clair \unsure{quén|qu'en} récollant $\Lambda_{xz}$ avec $\Lambda_{\unsure{zz|zy}}$ selon $\{z\}$ on obtient $\Lambda_{xy}\subset V$.

On a vu que $U$ et $V$ sont des connexes il sont necessairement ouverts car il sont les composantes connexes de $\mathbb{R}^{2}-\Gamma$, et $\Gamma$ est compact donc fermé dans $\mathbb{R}^{2}$\unsure{,|;|)}

\hfill \unsure{c.q.s.d.|c.q.f.d.}
% archive sheet number: 42
\scanpage{44}\mspage{XVI}
\uline{Lemme} $\underline{\overline{\mathrm{XVII}}}$. Si $x\in\Gamma$ et $D_x$ est un disque de centre $x$ qui ne coupe que les arêtes auxquelles $x$ appartient alors $D_x\cap U\ne\unsure{\phi|\varnothing}$, $D_x\cap V\ne\phi$.

\uline{Démonstration} : D'aprés la remarque qui suit la def. 1 il suffit d'exhiber un disque $D_x$ qui ne coupe que les arètes auxquelles $x$ appartient et tel que $D_x\cap U\ne\phi$, $D_x\cap V\ne\phi$

Si $x\in\Gamma^{\unsure{1|\prime}}_{ss'}-\{s,s'\}$ soit $D_x$ un disque de centre $x$ vérifiant : a) $D_x\subset\mathbb{R}^{2}-\Lambda^{2}$ ; b) $D_x$ satisfait b) de la \unsure{defnition|definition} 1 par rapport à $\Lambda^{1}$, donc $D_x\cap\mathrm{ext}\,\Lambda^{1}\ne\phi$ et par conséquent $D_x\cap V\ne\phi$. Aussi $D_x\subset\mathrm{ext}\,\Lambda^{2}$ et $D_x\cap\mathrm{int}\,\Lambda^{1}\ne\phi$ entraîne $D_x\cap V\ne\phi$. Si $x\in\Gamma^{2}_{ss'}-\{s,s'\}$ un argument analogue montre que $D_x\cap U\ne\phi$ et $D_x\cap V\ne\phi$.

Soit maintenant un disque $D_s$ qui satisfait b) de la \unsure{defnition|definition} 1 par rapport a $\Lambda^{1}$ et à $\Lambda^{2}$ on sait par le lemme \inserted{\textcircled{\scriptsize\unsure{XIV|XV}}} que
\[
\begin{aligned}
D_s\cap{}&\mathrm{ext}\,\Lambda^{1}\ne\phi
  &\quad\text{et}\quad D_s\cap{}&(\mathrm{int}\,\Lambda^{1}-\mathrm{int}\,\Lambda^{2})\\[-.5ex]
 &\hphantom{\mathrm{e}}\inserted{\begin{gathered}[t]\cap\\V\end{gathered}}
  &\begin{gathered}[t]\vphantom{\cap}\\\ne\phi\end{gathered}
  &\hphantom{(\mathrm{int}\,\Lambda^{1}-}\inserted{\begin{gathered}[t]\cap\\U\end{gathered}}
\end{aligned}
\]
ce qui suffit à prouver le lemme \unsure{s}pour $s$, un argument analogue démontre le lemme pour $s'$.

La \unsure{conjonction|conjonetion} du lemme $\underline{\overline{\mathrm{XVI}}}$ et du lemme $\underline{\overline{\mathrm{XVII}}}$ montre le \unsure{thēorēme|théorème} 1.

4- \uline{Preuve du \unsure{théorēme|théorème} 2} (On suppose que tout \inserted{polygone} $\Gamma$ est un $\unsure{J|j|y}$-polygone)

On procède par récurrence sur $\deleted{\unsure{s|S}}\mathop{\mathrm{card}} S(\Gamma)-2$.

On choisit $s,s'\in S(\Gamma)$ satisfaisant la condition du lemme XI, i.e. $s'$ est visible de $s$ de l'\unsure{interieur|intérieur}.

\deleted{\unsure{Posons}} \inserted{Soient} $\Gamma^{1}_{ss'}$ et $\Gamma^{2}_{ss'}$ \deleted{\unsure{égale|égals}} les \unsure{deux|dex} arcs de polygone de $\Gamma$ réliant $s$ à $s'$ décrits dans la Remarque 1, et soit $\Lambda_{ss'}$ l'arc de polygone $\Lambda_{ss'}=[s,s']$. Appelons $\Lambda^{1}$ \deleted{\unsure{l'arc|le}} (resp. $\Lambda^{2}$) le polygone \deleted{\unsure{f}}obtenu à partir de $\Gamma^{1}_{ss'}$ (resp. $\Gamma^{2}_{ss'}$) et $\Lambda_{ss'}$ par \unsure{rēcollement|récollement} selon $\{s,s'\}$.

Il est clair que $\inserted{\mathrm{card}}\,S(\Lambda^{1})<\mathop{\mathrm{card}} S(\Gamma)$ (resp. $\mathop{\mathrm{card}} S(\Lambda^{2})<\mathop{\mathrm{card}} S(\Gamma)$), autrement
% archive sheet number: 43
\scanpage{45}\mspage{XVII}
$[s,s']\unsure{\subset|\in} A(\Lambda^{1})$ (resp. $A(\Lambda^{2})$) ce qui est absurde \unsure{puìsque|puisque}
\[
{]s,s'[}\subset\mathbb{R}^{2}-\Lambda^{1}\quad(\text{resp. }\mathbb{R}^{2}-\Lambda^{2}).
\]

\uline{Lemme} $\underline{\overline{\mathrm{XVIII}}}$. $\Lambda^{2}-\{s,s'\}\deleted{\unsure{\supset}}\subset\mathrm{ext}\,\Lambda^{1}$ et $\Lambda^{1}-\{s,s'\}\subset\mathrm{ext}\,\Lambda^{2}$.

\uline{Démonstration} : Soit $D_s$ un disque de centre $s$ satisfaisant b) de la def. 1 par rapport a $\Lambda^{1}$ \deleted{et}, à $\Lambda^{2}$ et \unsure{\={a}|à} $\Lambda$.

Observons d'abord que le \unsure{secteur|seeteur} angulaire exterieur \inserted{\tikz[baseline=(n.base)]{\node[inner sep=2pt] (n) {$D_s\cap\mathrm{ext}\,\Lambda$}; \draw (n.west) .. controls (n.north west) and (n.north east) .. (n.east) .. controls (n.south east) and (n.south west) .. cycle;}} determiné par $\Lambda$ vérifie
\[
D_s\cap\mathrm{ext}\,\Lambda\subset D_s\cap\mathrm{ext}\,\Lambda^{1}\quad(\text{resp. }D_s\cap\mathrm{ext}\,\Lambda^{2}).
\]
En effet tout point $x$ de $D_s\cap\mathrm{ext}\,\Lambda$ peut être relié à un point $y$ de $\mathbb{R}^{2}$ tel que
\[
y\notin(\mathrm{int}\,\Lambda^{1}\cup\mathrm{int}\,\Lambda^{2}\cup\mathrm{int}\,\Lambda)\cup(\Lambda^{1}\cup\Lambda^{2}\cup\Lambda)
\]
\deleted{puisque} (on peut toujours trouver un tel $y$), par un arc de polygone $\Lambda_{xy}$ avec $\Lambda_{xy}\cap\Gamma=\unsure{\phi|\varnothing}$ \deleted{\unsure{et}} tel que $\Lambda_{xy}\subset\mathrm{ext}\,\Lambda$. Il est clair que $\Lambda_{xy}\cap\Lambda^{i}=\phi$ puisque $\Lambda_{xy}\cap{]s,s'[}=\phi$ car $ {]s,s'[}\deleted{\unsure{\not\subset}}\subset\mathrm{int}\,\Lambda^{i}\deleted{\unsure{\ne\phi|=\phi}}$, comme $y\notin\mathrm{int}\,\Lambda^{i}$, $i=1,2$, cela montre que $x\in\mathrm{ext}\,\Lambda^{i}$, $i=1,2$.

Cela permet de prouver que $(\Lambda^{2}\deleted{\unsure{\cap}}\mathbin{\unsure{\dot{-}|-}}{]s,s'[})\subset\overline{\mathrm{ext}\,\Lambda^{1}}$, mais le fait que $\Lambda^{2}-{]s,s'[}\subset\mathbb{R}^{2}-\Lambda^{1}$ entraîne que $\Lambda^{2}-{]s,s'[}\subset\mathrm{ext}\,\Lambda^{1}$. On montre aussi \inserted{avec le même argument} que $\Lambda^{1}-{]s,s'[}\subset\mathrm{ext}\,\Lambda^{1}$, c.q.f.d.

\uline{Lemme} $\underline{\overline{\mathrm{XIX}}}$. $\mathrm{int}\,\Lambda^{1}\subset\mathrm{ext}\,\Lambda^{2}$, $\mathrm{ext}\,\Lambda^{1}\supset\mathrm{int}\,\Lambda^{2}$.

\uline{Demonstration} : on montre que tout point de $\mathrm{int}\,\Lambda^{1}$ peut être relié\deleted{\unsure{r|e}} à un point de $\mathrm{ext}\,\Lambda^{2}$ par un arc de polygone qui ne coupe pas $\Lambda^{2}$. Comme $\Lambda^{\unsure{1|4|+}}-\{s,s'\}\subset\mathrm{ext}\,\Lambda^{2}$ soit $x\in\Lambda^{1}-\{s,s'\}$ alors \inserted{il \unsure{exìste|existe}} $D_x\subset\mathrm{ext}\,\Lambda^{2}$ satisfaisant b) de la def. 1. par rapport a $\Lambda^{1}$, donc $D_x\cap\mathrm{int}\,\Lambda^{1}\ne\phi$. Soi\inserted{en}t $z\in D_x\cap\mathrm{int}\,\Lambda^{1}$ et $y\in\mathrm{int}\,\Lambda^{1}$ il \unsure{exìste|existe} un arc de polygone $\Lambda_{zy}\subset\mathrm{int}\,\Lambda^{1}$ réliant $z\,(\in\mathrm{ext}\,\Lambda^{2})$ à $y$.
% archive sheet number: 44
\scanpage{46}
\figure{\begin{tikzpicture}[x=.021cm,y=-.021cm,line join=round,line cap=round]
% Upper drawing. Its lower vertex and the second drawing share this coordinate system.
\draw[black!78,line width=.85pt]
  (771,1187) .. controls (815,1219) and (864,1258) .. (908,1324)
  -- (1028,1205) -- (969,1235) -- (1029,1127) -- (1130,1253)
  -- (886,1496) .. controls (857,1397) and (809,1279) .. (771,1187);
\draw[black!63,line width=.55pt]
  (992,1240)--(1025,1207)--(978,1231);
\draw[black!46,line width=.55pt]
  (777,1187)--(795,1188) (812,1190)--(830,1190)
  (875,1195)--(892,1195) (936,1200)--(955,1201)
  (973,1203)--(991,1204) (1001,1204)--(1014,1205);
\draw[black!20,line width=.28pt]
  (908,1326) .. controls (906,1384) and (897,1455) .. (886,1496);
\draw[black!70,line width=1.1pt]
  (908,1328)--(909,1345) (907,1358)--(905,1377)
  (903,1392)--(903,1408) (899,1422)--(896,1443)
  (892,1464)--(887,1494);
\draw[black!48,thin] (886,1494)--(884,1512);
\node[anchor=south east,inner sep=1pt] at (765,1184) {$s'$};
\node[anchor=north west,inner sep=1pt] at (913,1324) {$s$};
\node[anchor=west,inner sep=0pt] at (1025,1210) {$s''$};
% Lower drawing.
\draw[black!60,line width=.8pt]
  (1028,1536) .. controls (1030,1602) and (1042,1690) .. (1055,1754)
  -- (1205,1535) -- (1076,1624) -- cycle;
\draw[line width=.7pt] (1028,1536)--(1076,1624)--(1089,1621);
\draw[black!45,line width=.42pt] (1080,1618)--(1205,1535);
\draw[black!27,line width=.35pt] (1030,1538)--(1038,1604);
\draw[black!48,line width=.42pt] (1028,1536)--(1036,1538);
% Isolated short ink stroke well below the four-cornered outline.
\draw[black!58,line width=.48pt]
  (1114,2020) .. controls (1115,2023) and (1113,2027) .. (1114,2032);
\end{tikzpicture}}

% margin note written sideways
% The writing is inverted (180 degrees) at the lower right.
\margin{ment\qquad $\deleted{\unsure{]}}[s,s']\cap$}
% archive sheet number: 45
\scanpage{47}\mspage{XVIII}
mais $\Lambda_{zy}\cap{]s,s'[}=\unsure{\phi|\varnothing}$ puisque $\Lambda_{zy}\subset\mathrm{int}\,\Lambda^1$, et par la même raison $\Lambda_{zy}\cap(\Lambda^2-\{s,s'\})=\phi$ car $\Lambda^2-\{s,s'\}\subset\mathrm{ext}\,\Lambda^1$ (cf. lemme $\overline{\underline{\mathrm{XVIII}}}$) cela prouve la \unsure{premiēre|première} observation, on obtient la seconde en changeant $\Lambda^1$ par $\Lambda^2$ dans la preuve ci-dessus, \unsure{c.q.f.d.|c.q.s.d.}

\uline{Lemme} $\overline{\underline{\mathrm{XX}}}$. Il y'a deux bons sommets de $\Gamma$, \inserted{non consécutifs si $\unsure{p|n}=$ \unsure{nb|nbre} de \unsure{côtés|cotés} de $\Gamma\geq4$,} \deleted{au moins} tels que les \inserted{\unsure{interieurs|intérieurs} des} triangles corespondants soient contenus dans $\mathrm{int}\,\Gamma$.

\uline{Démonstration} : L'\unsure{hypothēse|hypothèse} de récurrence appliquée à $\Lambda^1$ et $\Lambda^2$ montre qu'il'y a un \uline{bon sommet}\margin{\unsure{tel que}\unsure{…}} % slanted
$t_1$ dans $\Lambda^1$ et un bon sommet $t_2$ dans $\Lambda^2$. Si $t_1\ne s,s'$ alors $t_1$ est un bon sommet de $\Gamma$ puisque $\mathrm{int}\,\Lambda^1\subset\mathrm{int}\,\Lambda$. (tout point de $\mathrm{int}\,\Lambda^1$ peut être relié à un point de $\mathrm{int}\,\Lambda$, par \unsure{example|exemple} en prenant un disque $D_x$, $x\in\Gamma^{\unsure{1|\prime}}_{ss'}-\{s,s'\}$ vérifiant b) de la def.1 par rapport à $\Gamma$ et à $\Lambda^1$ on voit que $\mathrm{int}\,\Lambda^1\cap\mathrm{int}\,\Lambda\ne\phi$....etc)

Dans le cas contraire sans perte de généralité on peut poser $t_1=s$. Deux cas \margin{?}peuvent se produire $ {]\text{\textcircled{\scriptsize \(s''\)}},s'[}\cap\Gamma=\phi$ ou bien $ {]s'',s'[}\cap\Gamma\ne\phi$ et dans ce cas $ {]s'',s'[}\cap S(\Gamma)\ne\phi$ puisque $ {]s'',s'[}\not\subset\Gamma$.

\margin{\deleted{\unsure{faux}}} % slanted
Si $\deleted{\unsure{[}}{]s'',s'[}\cap\Gamma=\phi$ il est clair que $s''$ est visible de $s'$ par l'\unsure{interieur|intérieur} (car le triangle de sommets $s,s',s''\subset\mathrm{int}\,\Gamma$). \unsure{Considerons|Considérons} alors $\unsure{\Lambda^{1'}_{\mkern5mu\cdot}|\Lambda^{1!}|\Lambda^{11}_{\mkern5mu\cdot}}$ le polygone \margin{!}obtenu en récollant selon \deleted{\illegible{1}} $\{s',s''\}$ l'arc de polygone $[s',s'']$ avec le seul des arcs de polygone de $\Lambda^1$ réliant $s$ à $s'$ qui ne contient pas $[s,s']$. Il est clair que $S(\Lambda^{\unsure{1'|11}})\subset S(\Gamma)$ et que l'ensemble de toutes les arêtes de $\Lambda^{1'}$ sauf une, savoir $[s',s'']$ est contenu dans $A(\Gamma)$, \unsure{mais|Mais} nécessairement \margin{?}l'ensemble des arètes de $\Lambda^{1'}$ contenues dans $A(\Gamma)$ est de cardinal \unsure{inferieur|inférieur} \margin{\unsure{so what|So what}} % slanted
au cardinal de l'ensemble des arêtes de $\Lambda^1$ contenues dans $A(\Gamma)$.

Si $ {]s'',s'[}\cap S(\Gamma)\ne\phi$ on choisit $z\in S(\Gamma)\cap{]s'',s'[}$, il est clair que $z$ est visible \inserted{de $s=t$} par l'\unsure{interieur|intérieur} (puisque $ {]s,z[}\subset\overset{\circ}{\text{triangle}}$ de sommets $s,s',s''$). à partir du couple $s,z$ on peut appliquer la même construction que ci-dessus pour aboutir aussi à un polygone $\Lambda^{1'}$ tel que : a) $S(\Lambda^{1'})\subset S(\Gamma)$; b)
% archive sheet number: 46
\scanpage{48}\mspage{XIX}
tou\deleted{\unsure{s}}\unsure{tes|ts} les \deleted{\unsure{so|sa}}\unsure{.} arêtes de $\Lambda^{\unsure{1'|11}}$ sauf une (savoir, $[z,s]$) appartiennent à $\deleted{\unsure{A|R}}\,A(\Gamma)$ et le cardinal de vérifie $A(\Lambda^{1'})\cap A(\Gamma)<A(\Gamma)\cap A(\Lambda^1)$.

Comme il \unsure{n'ya|n'y a} qu'un nombre fini d'arêtes de $\Gamma$ il est clair qu'on \deleted{\unsure{bien}} a :

a) Il existe $n$ tel que $\deleted{\unsure{A|R}}\,\Lambda^{1(n)}$ (Notation definie par \unsure{récurrence|récurence} de \unsure{foçon|façon} \unsure{evidente|évidente}) contienne un \inserted{bon} sommet $t$ tel que les deux arêtes qui contiennent $t$ soient des arêtes $\Gamma$, et par conséquent $t$ est aussi un bon sommet de $\Gamma$ compte tenu du lemme $\overline{\underline{\mathrm{XIX}}}$, appliqué $n$-fois, \inserted{ou bien,}

b) il existe $n$ tel que $\mathrm{card}\,(A({\unsure{\Lambda|A}}^{1(n)})\cap A(\Gamma))=2$ donc $\mathrm{card}\,A^{1(n)}=3$ et il est clair que le \inserted{seul} sommet de $\Lambda^{1(n)}$ qui appartient aux deux arêtes de $\Lambda^{1(n)}$ contenues dans $A(\Gamma)$ est un bon sommet de $\Gamma$, compte tenu encore une fois du lemme $\overline{\underline{\mathrm{XIX}}}$\inserted{appliqué $n$ fois\unsure{,|.} Le même argument avec $\Lambda^2$ prouve le reste,}, \unsure{c.q.f.d.|c.q.s.d.}

\uline{Lemme $\overline{\underline{\mathrm{XXI}}}$} : Si $\Gamma$ est un polygone non-convexe il existe un bon sommet de $\Gamma$ tel que le \inserted{\unsure{interieur|intérieur} du} triangle corespondant est contenu dans $\mathrm{ext}\,\Gamma$.

\uline{Demonstration} : le lemme $\overline{\underline{\mathrm{VIII}}}$ montre \unsure{qu il|qu'il} \unsure{exìste|existe} un sommet $s$ de $\Gamma$ et un sommet $s'$ visible de $s$ par l'\unsure{exterieur|extérieur} de $\Gamma$. On conserve ici les notations du \unsure{numèro|numéro|numero} $3$. On considère $\Lambda^2$ (defini comme dans le numéro $3$) il est clair \margin{\unsure{faux}, \unsure{si on} \unsure{ne choisit} \unsure{pas bien} \unsure{ses|s.} \unsure{notations|solutions} \unsure{--|…}} % slanted
que $\mathrm{int}\,\Lambda^2\subset\mathrm{ext}\,\Gamma$ car $V=\mathrm{ext}\,\Lambda^1\cup\mathrm{int}\,\Lambda^2\cup{]s,s'[}$, et $V$ est non \deleted{\unsure{co|lo}}-relativement compact. \deleted{L'argument du}\inserted{Le} lemme \unsure{précedent|précédent} appliqué a $\Lambda^2$ montre que l'on peut trouver un \deleted{\unsure{trian|triang}} bon sommet $s$ de $\Lambda^2$ tel que : le triangle \unsure{corespondent|corespondant} soit contenu dans $\mathrm{ext}\,\Gamma$ et les deux arêtes de $\Lambda^2$ auxquelles \margin{\uline{\unsure{douteux|donc ?}}} % slanted
$s$ appartient soient des arêtes de $\Gamma$, ce qui permet de prouver le lemme.

\deleted{\unsure{Le|Les}}\unsure{Le|Lee} lemme $\overline{\underline{\mathrm{XX}}}$ le lemme $\overline{\underline{\mathrm{XXI}}}$ joints au fait que dans tout polygone convexe il 'y a au moins trois bons sommets, \unsure{montent|montrent} le théorème $2$, c.q.f.d.

\begin{flushright}
le 15 \unsure{juillet|jullet} 1976, Montpellier\\
\unsure{Mauré|Mouré|Maury|Maurice}.
\end{flushright}
% archive sheet number: 47
